\chapter*{APPENDICES}
\addcontentsline{toc}{chapter}{APPENDICES}

\appendix

% Appendix A on the same page
{\centering
	\textbf{Appendix A\\Step-by-Step Calculations}\par
	\addcontentsline{toc}{chapter}{Appendix A: Step-by-Step Calculation}
}

\section*{Detachable Power Source}
The battery capacity for the detachable power source is calculated based on expected daily consumption and system voltage.  
The battery capacity is calculated as:
\begin{align*}
	\text{Battery Capacity} &= \frac{\text{Total Daily Consumption}}{\text{System Voltage}} \\ 
	&= \frac{500 \text{ Wh}}{12 \text{ V}} \\ 
	&= 41.67 \text{ Ah} + 0.858 \text{ Ah} \\ 
	\text{Battery Capacity} &\approx 43 \text{ Ah}
\end{align*}
\noindent{Use:} 50Ah \ce{LiFePO4} Battery

\section*{Station}
The station battery is sized to accommodate the detachable battery and additional reserve.  
Let BCDPS = Battery Capacity of Detachable Power Source. 

The station battery is sized as:
\begin{align*} 
	\text{Battery Capacity} &= (\text{BCDPS} + 35\%) + 0.858\text{Ah} \\
	&= 68.358\text{Ah} \\
	\text{Battery Capacity} &\approx 70 \text{Ah}
\end{align*}
\noindent{Use:} 70Ah \ce{LiFePO4} Battery.

\section*{Panel Sizing}
Solar panel power is calculated to meet the daily energy requirements within average peak sun hours.  
\begin{align*}
	\text{PV Power} &= \frac{\text{Total Daily Consumption}}{\text{Sun Peak Hours}} \\
	&= \frac{500\text{Wh}}{3.5\text{h}} \\
	&= 142.857\text{W} \\
	\text{PV Power} &= 200\text{W}
\end{align*}
Use: 200W Solar Panel.

\section*{Inverter Sizing}
The inverter is sized to handle the system load with a safety margin.  
\begin{align*}
	\text{Inverter Size} &= 250W + 20\% \text{ allowance} \\
	\text{Inverter Size} &= 300W
\end{align*}
Use: 300W pure sine wave inverter.

\section*{Solar Charge Controller Sizing}
The charge controller current is selected based on PV array power and system voltage.  
\begin{align*}
	I_{cc} &= \frac{\text{Total PV Power}}{\text{System Voltage}} \\
	&= \frac{200W}{12V} \\
	&= 16.67A \\
	I_{cc} &= 20A
\end{align*}
Use: 1 pc — MPPT 12V or 24V Auto Adapt 20A.

\section*{Breaker Sizing}
Circuit breakers are sized to protect each part of the system.  
\begin{enumerate}[label=\alph*.]
	\item Panel - SCC
	\begin{align*}
		\text{Circuit Breaker} &= 11.85A \times \text{Over Radiance Factor} \times \text{Safety Factor} \\
		&= 18.5A \\
		\text{Circuit Breaker} &= 20A
	\end{align*}
	
	\item SCC - Battery
	\begin{align*}
		\text{Circuit Breaker} &= \frac{\text{Total Array Wattage}}{\text{Battery Bank Nominal Voltage}} \times \text{Safety Factor} \\
		&= \frac{200W}{12V} \times 1.25 \\
		\text{Circuit Breaker} &=20.8A
	\end{align*}
	
	\item Battery - Power Inverter
	\begin{align*}
		\text{Circuit Breaker} &= \frac{\text{Power Inverter}}{\text{Power Bank Nominal Voltage}} \\
		&= \frac{1000W}{12V}\\
		&= 83.33A\\
		\text{Circuit Breaker} &= 80A
	\end{align*}
	
	\item Inverter - Load (AC Breaker)\\
	BBNV = Battery Bank Nominal Voltage
	\begin{align*}
		\text{Circuit Breaker} &= \frac{\text{Safety Load}}{\text{Inverter Voltage Output}} \times \text{BBNV} \\
		&= \frac{700W}{220V} \times 1.25V \\
		\text{Circuit Breaker} &= 3.98A
	\end{align*}
\end{enumerate}

\section*{Wire Sizing}
Wire sizes are determined to handle current and limit voltage drop.  
\begin{enumerate}[label=\alph*.]
	\item Panel - SCC
	\begin{align*}
		\text{VDI} &= \frac{10.85 \times 12}{18.5 \times 2} \\
		&= 3.53\\
		\text{Wire Size} &= 10AWG
	\end{align*}
	\item SCC - Battery
	\begin{align*}
		\text{VDI} &= \frac{20 \times 4}{12.8 \times 2} \\
		&= 3.125\\
		\text{Wire Size} &= 10AWG
	\end{align*}
\end{enumerate}

\nextbreak
\section*{Calculation of Appliances Wattage Usage}

This section presents the computation of power consumption for commonly used household and electronic appliances. The energy usage of each device is calculated using the formula:

\begin{equation}
	E = P \times t
\end{equation}

where:  
\begin{itemize}
	\item $E$ = Energy consumed (in watt-hours, Wh)
	\item $P$ = Power rating of the appliance (in watts, W)
	\item $t$ = Usage time (in hours, h)
\end{itemize}

For example, a laptop rated at 65~W that operates for 4~hours per day will consume:
\[
E = 65~\text{W} \times 4~\text{h} = 260~\text{Wh/day}.
\]

\begin{table}[H]
	\centering
	\caption{Estimated Wattage Usage of Common Appliances}
	\label{tab:appliance_wattage}
	\begin{tabular}{|p{3cm}|c|c|c|c|}
		\hline
		\textbf{Appliance} & \textbf{Power Rating (W)} & \textbf{Usage Time (h/day)} & \textbf{Energy (Wh/day)}\\ \hline
		Smartphone (charging) & 15 & 2 & 30 \\ \hline
		Laptop & 65 & 4 & 260 \\ \hline
		Electric Fan & 75 & 8 & 600 \\ \hline
		Rice Cooker & 700 & 1 & 700 \\ \hline
		Refrigerator (inverter type) & 150 & 24 & 3600 \\ \hline
		\textbf{Total Daily Energy Usage} & --- & --- & \textbf{5190 Wh/day} \\ \hline
	\end{tabular}
\end{table}

The total estimated daily consumption of the listed appliances is approximately \textbf{5.19~kWh/day}. This serves as the baseline for sizing the solar-powered system and determining the battery capacity required for sustained operation during power interruptions.


